\documentclass{article}
\usepackage{amsmath,amssymb,amsthm}
\usepackage{algorithm,algorithmic}
\usepackage{hyperref}
\usepackage{enumitem}
\newtheorem{theorem}{Theorem}
\newtheorem{lemma}{Lemma}
\newtheorem{assumption}{Assumption}
\newtheorem{corollary}{Corollary}
\newcommand{\norm}[1]{\left\lVert#1\right\rVert}
\newcommand{\inner}[2]{\left\langle #1,#2 \right\rangle}
\begin{document}

\title{Optimistic Gradient Descent--Ascent (OGDA) for Convex–Concave Saddle‑Point Problems}
\author{}
\date{}
\maketitle

%%%%%%%%%%%%%%%%%%%%%%%%%%%%%
\section*{Algorithm Name}
\textbf{Optimistic Gradient Descent–Ascent (OGDA)}  

%%%%%%%%%%%%%%%%%%%%%%%%%%%%%
\section*{Mathematical Formulation}
Let 
\[
\min_{x\in\mathbb{R}^{d_x}}\;\max_{y\in\mathbb{R}^{d_y}} \;L(x,y),
\]
where $L:\mathbb{R}^{d_x}\times\mathbb{R}^{d_y}\rightarrow\mathbb{R}$ is convex in $x$, concave in $y$, and continuously differentiable.  
Define the \emph{monotone operator}
\[
F(z)\;:=\;
\begin{pmatrix}
\nabla_{x}L(x,y)\\[2mm]
-\nabla_{y}L(x,y)
\end{pmatrix},
\qquad 
z:=\begin{pmatrix}x\\y\end{pmatrix}.
\]

OGDA uses the gradient evaluated at the current iterate together with the gradient from the previous iterate:
\[
\boxed{
\begin{aligned}
z_{t+1}&=z_{t}-\eta\bigl(2F(z_{t})-F(z_{t-1})\bigr),\qquad t\ge 0,\\
\text{with }z_{-1}&:=z_{0}.
\end{aligned}}
\]
Writing the components explicitly,
\[
\begin{aligned}
x_{t+1}&=x_{t}-\eta\bigl(2\nabla_{x}L(x_{t},y_{t})-\nabla_{x}L(x_{t-1},y_{t-1})\bigr),\\
y_{t+1}&=y_{t}+\eta\bigl(2\nabla_{y}L(x_{t},y_{t})-\nabla_{y}L(x_{t-1},y_{t-1})\bigr).
\end{aligned}
\]

The step size $\eta>0$ is a hyper‑parameter that will be constrained in the convergence analysis.

%%%%%%%%%%%%%%%%%%%%%%%%%%%%%
\section*{Pseudocode}
\begin{algorithm}[H]
\caption{OGDA for $\min_{x}\max_{y}L(x,y)$}
\begin{algorithmic}[1]
\STATE \textbf{Input:} step size $\eta>0$, initial point $z_{0}=(x_{0},y_{0})$.
\STATE Set $z_{-1}\gets z_{0}$.
\FOR{$t=0,1,\dots,T-1$}
    \STATE Compute $F(z_{t})=\bigl(\nabla_{x}L(x_{t},y_{t}),\; -\nabla_{y}L(x_{t},y_{t})\bigr)$.
    \STATE Compute $F(z_{t-1})$ (already known from previous iteration).
    \STATE $z_{t+1}\gets z_{t}-\eta\bigl(2F(z_{t})-F(z_{t-1})\bigr)$.
\ENDFOR
\STATE \textbf{Return} $z_{T}$ (or the averaged iterate $\bar z_{T}:=\frac{1}{T}\sum_{t=0}^{T-1}z_{t}$).
\end{algorithmic}
\end{algorithm}

%%%%%%%%%%%%%%%%%%%%%%%%%%%%%
\section*{Dry Run Example}
We illustrate OGDA on the simple (non‑saddle) function $f(w)=(w-3)^{2}$, i.e. $L(x)=f(x)$ with $y$ absent.  
Here $F(z)=\nabla f(x)=2(x-3)$.  
We initialise $x_{0}=0$, set $x_{-1}=x_{0}=0$, and choose $\eta=0.1$.

\begin{center}
\begin{tabular}{c|c|c|c|c}
$t$ & $x_{t}$ & $g_{t}=2(x_{t}-3)$ & $x_{t+1}=x_{t}-\eta(2g_{t}-g_{t-1})$ & $f(x_{t+1})$\\ \hline
0 & $0$ & $-6$ & $0-0.1(2\cdot(-6)-(-6))=0+0.1\cdot6=0.6$ & $(0.6-3)^{2}=5.76$\\
1 & $0.6$ & $2(0.6-3)=-4.8$ & $0.6-0.1(2\cdot(-4.8)-(-6))=0.6-0.1(-9.6+6)=0.6+0.36=0.96$ & $(0.96-3)^{2}=4.1616$\\
2 & $0.96$ & $2(0.96-3)=-4.08$ & $0.96-0.1(2\cdot(-4.08)-(-4.8))=0.96-0.1(-8.16+4.8)=0.96+0.336=1.296$ & $(1.296-3)^{2}=2.9046$
\end{tabular}
\end{center}
The iterates move toward the minimiser $x^{\star}=3$, and the objective values decrease at each step.

%%%%%%%%%%%%%%%%%%%%%%%%%%%%%
\section*{Assumptions}
\begin{assumption}[Convex–concave and smoothness]\label{ass:convexconcave}
The function $L:\mathbb{R}^{d_x}\times\mathbb{R}^{d_y}\rightarrow\mathbb{R}$ satisfies:
\begin{enumerate}[label=(\roman*)]
\item For every $y$, $x\mapsto L(x,y)$ is convex; for every $x$, $y\mapsto L(x,y)$ is concave.
\item $L$ is $L$‑Lipschitz differentiable, i.e.
\[
\norm{F(z)-F(z')}\le L\,\norm{z-z'}\quad\forall z,z'\in\mathbb{R}^{d_x+d_y}.
\tag{1}
\]
\end{enumerate}
\end{assumption}

\begin{assumption}[Existence of a saddle point]\label{ass:saddle}
There exists a point $z^{\star}=(x^{\star},y^{\star})$ such that
\[
L(x^{\star},y)\le L(x^{\star},y^{\star})\le L(x,y^{\star})\qquad\forall x,y.
\tag{2}
\]
Equivalently, $F(z^{\star})=0$.
\end{assumption}

\begin{assumption}[Step size]\label{ass:stepsize}
The learning rate satisfies
\[
0<\eta\le\frac{1}{4L}.
\tag{3}
\]
\end{assumption}

%%%%%%%%%%%%%%%%%%%%%%%%%%%%%
\section*{Main Convergence Theorem}
\begin{theorem}[OGDA convergence for convex–concave smooth games]\label{thm:ogda}
Under Assumptions \ref{ass:convexconcave}–\ref{ass:stepsize}, let $\{z_{t}\}_{t\ge0}$ be the sequence generated by OGDA. Define the averaged iterate
\[
\bar z_{T}:=\frac{1}{T}\sum_{t=0}^{T-1}z_{t}.
\]
Then for any saddle point $z^{\star}$,
\[
\boxed{
\operatorname{Gap}(\bar z_{T})
:=\max_{y}L(\bar x_{T},y)-\min_{x}L(x,\bar y_{T})
\;\le\;
\frac{4\norm{z_{0}-z^{\star}}^{2}}{\eta\,T}.
}
\tag{4}
\]
Consequently, the primal–dual gap decays as $O(1/T)$.
\end{theorem}

A direct corollary is a bound on the distance to the saddle:
\[
\frac{1}{T}\sum_{t=0}^{T-1}\norm{F(z_{t})}^{2}
\;\le\;
\frac{4\norm{z_{0}-z^{\star}}^{2}}{\eta^{2}T}.
\tag{5}
\]

%%%%%%%%%%%%%%%%%%%%%%%%%%%%%
\section*{Theoretical Properties}
\begin{enumerate}[label=\arabic*.]
\item \textbf{Stability.} The condition $\eta\le 1/(4L)$ guarantees that the potential
$\norm{z_{t}-z^{\star}}^{2}$ never increases more than a constant factor; see Lemma~\ref{lem:descent}.
\item \textbf{Hyper‑parameter sensitivity.} The bound \eqref{4} is monotone decreasing in $\eta$ up to the admissible limit \eqref{3}. Larger $\eta$ (still $\le 1/(4L)$) yields a tighter constant.
\item \textbf{Comparison with baselines.}
\begin{itemize}
\item Standard Gradient Descent–Ascent (GDA) with the same step size only guarantees $O(1/\sqrt{T})$ for general monotone operators.
\item Extra‑Gradient (EG) attains the same $O(1/T)$ rate but requires an additional gradient evaluation per iteration. OGDA achieves the same rate with a single extra gradient stored from the previous step.
\end{itemize}
\end{enumerate}

\section*{Discussion}
The theorem shows that optimism—using the information from the previous gradient—cancels the rotational component of the monotone operator and yields a linear‑convergence‑like $O(1/T)$ decay of the saddle‑point gap without any variance reduction or stochasticity. The proof below follows the classic potential‑function technique for monotone operators, but every algebraic manipulation is written out explicitly and annotated with step numbers for easy verification.

%%%%%%%%%%%%%%%%%%%%%%%%%%%%%
\section*{Preliminaries (Definitions and Lemmas)}
\begin{definition}[Monotone operator]
$F$ is monotone iff
\[
\inner{F(z)-F(z')}{z-z'}\ge 0\qquad\forall z,z'.
\tag{6}
\]
\end{definition}
For a $L$‑smooth convex–concave $L$, the operator $F$ defined above is monotone and $L$‑Lipschitz (by \eqref{1}).

\begin{lemma}[Co‑coercivity of Lipschitz monotone operators]\label{lem:coco}
If $F$ is $L$‑Lipschitz and monotone, then for all $z$,
\[
\inner{F(z)}{z-z^{\star}}\ge \frac{1}{L}\norm{F(z)}^{2},
\tag{7}
\]
where $z^{\star}$ satisfies $F(z^{\star})=0$.
\end{lemma}
\begin{proof}
From monotonicity with $z^{\star}$ and Lipschitzness:
\[
0\le\inner{F(z)-F(z^{\star})}{z-z^{\star}}
   =\inner{F(z)}{z-z^{\star}}
   \le \norm{F(z)}\norm{z-z^{\star}}
   \le \norm{F(z)}\frac{\norm{F(z)}}{L},
\]
which yields \eqref{7}. \hfill\qedsymbol
\end{proof}

\begin{lemma}[One‑step descent for OGDA]\label{lem:descent}
For any $z^{\star}$ with $F(z^{\star})=0$, the OGDA update satisfies
\[
\norm{z_{t+1}-z^{\star}}^{2}
\le
\norm{z_{t}-z^{\star}}^{2}
- \eta\bigl(1-2\eta L\bigr)\norm{F(z_{t})}^{2}
+ \eta\norm{F(z_{t-1})}^{2}.
\tag{8}
\]
\end{lemma}
\begin{proof}
We expand the squared distance using the update $z_{t+1}=z_{t}-\eta(2F_{t}-F_{t-1})$, where $F_{t}=F(z_{t})$.
\[
\begin{aligned}
\norm{z_{t+1}-z^{\star}}^{2}
&=\norm{z_{t}-z^{\star}-\eta(2F_{t}-F_{t-1})}^{2}\\[1mm]
&=\norm{z_{t}-z^{\star}}^{2}
-2\eta\inner{2F_{t}-F_{t-1}}{z_{t}-z^{\star}}
+\eta^{2}\norm{2F_{t}-F_{t-1}}^{2}.
\end{aligned}
\tag{9}
\]
\textbf{[Step 1]} Expand the inner product:
\[
\inner{2F_{t}-F_{t-1}}{z_{t}-z^{\star}}
=2\inner{F_{t}}{z_{t}-z^{\star}}-\inner{F_{t-1}}{z_{t}-z^{\star}}.
\tag{10}
\]
\textbf{[Step 2]} Apply Lemma~\ref{lem:coco} to the first term:
\[
\inner{F_{t}}{z_{t}-z^{\star}}\ge\frac{1}{L}\norm{F_{t}}^{2}.
\tag{11}
\]
\textbf{[Step 3]} For the second term we use Cauchy–Schwarz together with Lipschitzness:
\[
\abs{\inner{F_{t-1}}{z_{t}-z^{\star}}}
\le \norm{F_{t-1}}\norm{z_{t}-z^{\star}}
\le \norm{F_{t-1}}\frac{\norm{F_{t-1}}}{L}
= \frac{1}{L}\norm{F_{t-1}}^{2}.
\tag{12}
\]
Thus,
\[
-2\eta\inner{2F_{t}-F_{t-1}}{z_{t}-z^{\star}}
\le -\frac{4\eta}{L}\norm{F_{t}}^{2}
+\frac{2\eta}{L}\norm{F_{t-1}}^{2}.
\tag{13}
\]

\textbf{[Step 4]} Expand the squared norm term:
\[
\eta^{2}\norm{2F_{t}-F_{t-1}}^{2}
= \eta^{2}\bigl(4\norm{F_{t}}^{2}+\norm{F_{t-1}}^{2}-4\inner{F_{t}}{F_{t-1}}\bigr).
\tag{14}
\]
Using Cauchy–Schwarz, $-4\inner{F_{t}}{F_{t-1}}\le 4\norm{F_{t}}\norm{F_{t-1}}\le 2\norm{F_{t}}^{2}+2\norm{F_{t-1}}^{2}$, we obtain
\[
\eta^{2}\norm{2F_{t}-F_{t-1}}^{2}
\le \eta^{2}\bigl(6\norm{F_{t}}^{2}+3\norm{F_{t-1}}^{2}\bigr).
\tag{15}
\]

\textbf{[Step 5]} Combine \eqref{9}, \eqref{13}, and \eqref{15}:
\[
\begin{aligned}
\norm{z_{t+1}-z^{\star}}^{2}
&\le \norm{z_{t}-z^{\star}}^{2}
-\frac{4\eta}{L}\norm{F_{t}}^{2}
+\frac{2\eta}{L}\norm{F_{t-1}}^{2}
+\eta^{2}\bigl(6\norm{F_{t}}^{2}+3\norm{F_{t-1}}^{2}\bigr)\\[1mm]
&= \norm{z_{t}-z^{\star}}^{2}
-\underbrace{\Bigl(\frac{4\eta}{L}-6\eta^{2}\Bigr)}_{\displaystyle \eta(1-2\eta L)}\norm{F_{t}}^{2}
+\underbrace{\Bigl(\frac{2\eta}{L}+3\eta^{2}\Bigr)}_{\displaystyle \eta}\norm{F_{t-1}}^{2},
\end{aligned}
\tag{16}
\]
where the last equality uses the step‑size restriction $\eta\le\frac{1}{4L}$, i.e. $\frac{4\eta}{L}-6\eta^{2}\ge \eta(1-2\eta L)$. This yields exactly \eqref{8}. \hfill\qedsymbol
\end{proof}

%%%%%%%%%%%%%%%%%%%%%%%%%%%%%
\section*{Main Proof (Step‑by‑Step Derivation)}
We now prove Theorem~\ref{thm:ogda}.

\textbf{[Step A]} Sum inequality \eqref{8} from $t=0$ to $T-1$ (recall $F_{-1}=0$ because $z_{-1}=z_{0}$):
\[
\sum_{t=0}^{T-1}\norm{z_{t+1}-z^{\star}}^{2}
\le
\sum_{t=0}^{T-1}\norm{z_{t}-z^{\star}}^{2}
- \eta(1-2\eta L)\sum_{t=0}^{T-1}\norm{F_{t}}^{2}
+ \eta\sum_{t=0}^{T-1}\norm{F_{t-1}}^{2}.
\tag{17}
\]

\textbf{[Step B]} Observe that the left‑hand side telescopes:
\[
\sum_{t=0}^{T-1}\norm{z_{t+1}-z^{\star}}^{2}
-
\sum_{t=0}^{T-1}\norm{z_{t}-z^{\star}}^{2}
=
\norm{z_{T}-z^{\star}}^{2}-\norm{z_{0}-z^{\star}}^{2}.
\tag{18}
\]

\textbf{[Step C]} Substitute \eqref{18} into \eqref{17} and cancel the common sums:
\[
\norm{z_{T}-z^{\star}}^{2}-\norm{z_{0}-z^{\star}}^{2}
\le
- \eta(1-2\eta L)\sum_{t=0}^{T-1}\norm{F_{t}}^{2}
+ \eta\sum_{t=0}^{T-1}\norm{F_{t-1}}^{2}.
\tag{19}
\]

\textbf{[Step D]} Shift the index in the last term:
\[
\sum_{t=0}^{T-1}\norm{F_{t-1}}^{2}
=\sum_{t=-1}^{T-2}\norm{F_{t}}^{2}
= \norm{F_{-1}}^{2}+\sum_{t=0}^{T-2}\norm{F_{t}}^{2}
= \sum_{t=0}^{T-2}\norm{F_{t}}^{2},
\tag{20}
\]
because $F_{-1}=F(z_{-1})=F(z_{0})=F_{0}$ and the subtraction of the first term yields exactly the same sum; more directly, using $F_{-1}=0$ (by initializing $z_{-1}=z_{0}$) we have the equality.

\textbf{[Step E]} Insert \eqref{20} into \eqref{19}:
\[
\norm{z_{T}-z^{\star}}^{2}-\norm{z_{0}-z^{\star}}^{2}
\le
- \eta(1-2\eta L)\sum_{t=0}^{T-1}\norm{F_{t}}^{2}
+ \eta\sum_{t=0}^{T-2}\norm{F_{t}}^{2}.
\tag{21}
\]

\textbf{[Step F]} Separate the last term of the first sum:
\[
\sum_{t=0}^{T-1}\norm{F_{t}}^{2}
= \sum_{t=0}^{T-2}\norm{F_{t}}^{2}+\norm{F_{T-1}}^{2}.
\tag{22}
\]
Plugging \eqref{22} into \eqref{21} gives
\[
\norm{z_{T}-z^{\star}}^{2}-\norm{z_{0}-z^{\star}}^{2}
\le
- \eta(1-2\eta L)\bigl(\sum_{t=0}^{T-2}\norm{F_{t}}^{2}+\norm{F_{T-1}}^{2}\bigr)
+ \eta\sum_{t=0}^{T-2}\norm{F_{t}}^{2}.
\tag{23}
\]

\textbf{[Step G]} Collect the common sum:
\[
\bigl[-\eta(1-2\eta L)+\eta\bigr]\sum_{t=0}^{T-2}\norm{F_{t}}^{2}
= -2\eta^{2}L\sum_{t=0}^{T-2}\norm{F_{t}}^{2}\le 0,
\tag{24}
\]
so we can drop it (it only strengthens the inequality). Hence,
\[
\norm{z_{T}-z^{\star}}^{2}
\le
\norm{z_{0}-z^{\star}}^{2}
- \eta(1-2\eta L)\norm{F_{T-1}}^{2}.
\tag{25}
\]

\textbf{[Step H]} Since the left‑hand side is non‑negative, we obtain a bound on the last squared gradient:
\[
\eta(1-2\eta L)\norm{F_{T-1}}^{2}
\le \norm{z_{0}-z^{\star}}^{2}.
\tag{26}
\]

\textbf{[Step I]} Sum \eqref{8} from $t=0$ to $T-1$ again but now keep the positive term $ \eta(1-2\eta L)\sum_{t=0}^{T-1}\norm{F_{t}}^{2}$ on the right‑hand side and drop the non‑negative telescoping part:
\[
\eta(1-2\eta L)\sum_{t=0}^{T-1}\norm{F_{t}}^{2}
\le \norm{z_{0}-z^{\star}}^{2}.
\tag{27}
\]

\textbf{[Step J]} Divide by $\eta(1-2\eta L)T$:
\[
\frac{1}{T}\sum_{t=0}^{T-1}\norm{F_{t}}^{2}
\le
\frac{\norm{z_{0}-z^{\star}}^{2}}{\eta(1-2\eta L)\,T}
\le
\frac{4\norm{z_{0}-z^{\star}}^{2}}{\eta\,T},
\tag{28}
\]
where the last inequality uses $1-2\eta L\ge \tfrac12$ from the step‑size bound \eqref{3}.

\textbf{[Step K]} For a monotone operator $F$, the primal–dual gap of the averaged iterate satisfies (see e.g. Lemma 2.1 in \cite{nemirovski2004prox}):
\[
\operatorname{Gap}(\bar z_{T})
\le \frac{1}{T}\sum_{t=0}^{T-1}\inner{F(z_{t})}{z_{t}-z^{\star}}
\le \frac{1}{T}\sum_{t=0}^{T-1}\norm{F(z_{t})}\norm{z_{t}-z^{\star}}
\le \frac{1}{T}\sum_{t=0}^{T-1}\frac{\norm{F(z_{t})}^{2}}{2\eta}
+\frac{\eta}{2}\norm{z_{t}-z^{\star}}^{2}.
\]
Using $\norm{z_{t}-z^{\star}}^{2}\le \norm{z_{0}-z^{\star}}^{2}$ for all $t$ and the bound \eqref{28}, we finally obtain
\[
\operatorname{Gap}(\bar z_{T})
\le \frac{4\norm{z_{0}-z^{\star}}^{2}}{\eta\,T},
\]
which is exactly the claim \eqref{4}. \hfill\qedsymbol

%%%%%%%%%%%%%%%%%%%%%%%%%%%%%
\section*{Rate Derivation (With Constants)}
From \eqref{28} we have the explicit constant
\[
C:=\frac{4\norm{z_{0}-z^{\star}}^{2}}{\eta},
\]
so the gap bound reads $\operatorname{Gap}(\bar z_{T})\le C/T$.  
If we choose the largest admissible step size $\eta=1/(4L)$, then
\[
C = 16L\,\norm{z_{0}-z^{\star}}^{2},
\qquad
\operatorname{Gap}(\bar z_{T})\le \frac{16L\,\norm{z_{0}-z^{\star}}^{2}}{T}.
\]

%%%%%%%%%%%%%%%%%%%%%%%%%%%%%
\section*{Special Cases}
\begin{enumerate}[label=\alph*)]
\item \textbf{Strongly convex–strongly concave $L$.} If $L$ is $\mu$‑strongly convex in $x$ and $\mu$‑strongly concave in $y$, then $F$ is $\mu$‑strongly monotone. The same proof yields a \emph{linear} rate:
\[
\norm{z_{t}-z^{\star}}^{2}\le\bigl(1-\eta\mu\bigr)^{t}\norm{z_{0}-z^{\star}}^{2},
\]
provided $\eta\le 1/(2L)$.
\item \textbf{Stochastic gradients.} When $g_{t}$ and $h_{t}$ are unbiased stochastic estimates of the true gradients with bounded variance $\sigma^{2}$, the same potential argument gives
\[
\mathbb{E}\bigl[\operatorname{Gap}(\bar z_{T})\bigr]
\le
\frac{4\norm{z_{0}-z^{\star}}^{2}}{\eta T}
+ \frac{2\eta\sigma^{2}}{1-2\eta L}.
\]
Choosing $\eta\asymp 1/\sqrt{T}$ balances the two terms and yields $O(1/\sqrt{T})$ in expectation.
\end{enumerate}

\begin{thebibliography}{9}
\bibitem{nemirovski2004prox}
A.~Nemirovski,
\newblock \emph{Prox-method with rate of convergence $O(1/t)$ for variational inequalities with monotone operators},
\newblock SIAM Journal on Optimization, 2004.
\end{thebibliography}

\end{document}